% Source type: Mahbubur Rahman, Chapter 5 PDF p.45 (constants).
\booksubheading{ধ্রুবক হিসেবে ব্যবহৃত নাম}

কোনো প্রোগ্রামে যে মান অপরিবর্তিত রাখার পরিকল্পনা করা হয়, তাকে ধ্রুবক
হিসেবে ব্যবহার করা যায়। Python-এ এমন নাম সাধারণত বড় হাতের অক্ষরে লেখা হয়,
যেমন \py{VAT\_RATE}। এটি নামকরণের রীতি; বড় হাতের নাম লিখলেই নতুন মান
assign করা বন্ধ হয়ে যায় না। তাই প্রোগ্রামারকে মানটি অপরিবর্তিত রাখতে হয়।

\textbf{উদাহরণ: নির্ধারিত হারে করসহ মূল্য।}
\begin{pythoncode}
VAT_RATE = 0.05
price = 800
vat = price * VAT_RATE
payable = price + vat
print(f"VAT = {vat:.2f}")
print(f"Payable = {payable:.2f}")
\end{pythoncode}
\textbf{Output:} \py{VAT = 40.00}, \py{Payable = 840.00}।
এখানে মূল দাম পরিবর্তিত হতে পারে; একই প্রোগ্রামের হিসাবগুলোতে করের হার
একই রাখতে \py{VAT\_RATE} ব্যবহার করা হয়েছে।

